\documentclass{article}
\usepackage[top=1in, bottom=1in, left=1.3in, right=1.3in]{geometry}
\usepackage{graphicx}
\usepackage{tcolorbox}
\tcbuselibrary{listings, breakable}
\usepackage{booktabs}
\usepackage{amsmath}
\usepackage{float}
\usepackage{amssymb}
\usepackage{xcolor}
\usepackage{hyperref}
\usepackage{caption}
\usepackage{subcaption}

% Inline flags for results that are pending confirmation/rerun, so nothing
% provisional silently reads as final.
%   \pending{...}     -- data that needs a rerun before the number is final
%   \todoconfirm{...} -- a scope question for Xuanhao/Dr. Liu, not yet answered
\newcommand{\pending}[1]{\textcolor{orange!70!black}{[PENDING RERUN: #1]}}
\newcommand{\todoconfirm}[1]{\textcolor{red!70!black}{[CONFIRM WITH XUANHAO: #1]}}

\title{\textbf{Downlink-Constrained Task Scheduling for Orbital Edge Computing}}
\author{Hemanth Sudhaharan \\ NICE Lab, North Carolina State University}
\date{\today}

\begin{document}
\maketitle

\begin{abstract}
This document works through the on-board compression and downlink
scheduling problem for Orbital Edge Computing (OEC): a residual-quantized
VAE (RQ-VAE) compresses each image at one of several selectable depths,
and a rolling-horizon MPC scheduler decides depth, timing, and routing
under time-varying satellite/ground-station connectivity. It covers the
system model and full scheduling formulation -- extending the original
task model and MPC objective with a unified utility function, a
congestion-aware routing variant, and hierarchical/two-level MPC couplings
-- and the experimental setup (constellation, link-capacity regimes, task
workload) those schedulers are evaluated against. Results, related work,
and a conclusion are being written up separately as the corresponding
experiments are finished.
\end{abstract}

\section{Introduction}
\label{sec:intro}

A satellite in low Earth orbit (LEO) can observe far more of the ground
than it can ever report. A single imaging pass over a wildfire, a flood,
or an agricultural region can generate far more raw pixels than the
satellite's downlink budget can move to a ground station before the next
useful decision has to be made — and a LEO satellite is only in radio
contact with any one ground station for a few minutes at a time before it
drops below the horizon. Someone, or something, has to decide which
images matter most, how aggressively to compress each one, and when to
send it. \emph{Orbital Edge Computing} (OEC) is the idea of making some of
that decision on board the satellite itself — compressing and
prioritizing at the edge of the network, rather than shipping raw data
down and sorting it out on the ground.

This paper studies that decision problem end to end: an on-board
residual-quantized variational autoencoder (RQ-VAE) compresses each image
at one of several selectable \emph{depths} (more refinement passes, larger
payload, higher fidelity), and a rolling-horizon Model Predictive Control
(MPC) scheduler decides, for every task competing for downlink capacity,
which depth to use and when and along which route to transmit it, under a
constellation topology and ground-station contact schedule that changes
every scheduling slot.

\paragraph{The problem with grading a scheduler on how the image looks.}
The natural way to score a compression-and-scheduling policy is to ask
whether the reconstructed image still \emph{looks} like the original —
standard fidelity metrics such as PSNR and LPIPS answer exactly that
question. But for a mission whose actual goal is to extract information
from the imagery (what land-cover class is this, where is the flood
boundary), \emph{looking similar} and \emph{remaining useful for the task}
are not the same property, and this paper's first and most consequential
finding is that they diverge sharply here: pixel-fidelity quality changes
by only $12\%$ across the compression-depth range we test, while the
imagery's usefulness for a real downstream semantic-segmentation task
changes by $247\%$ over the identical range (\S\ref{sec:downstream}). A
scheduler rewarded on the flat proxy has almost nothing to gain by
reasoning carefully about depth; a scheduler rewarded on the true
task-relevant signal has a great deal to gain. This single substitution —
replacing a pixel-fidelity proxy with a measured downstream-task quality
signal — changes MPC's measured advantage over the best fixed-depth
baseline from a marginal $1.2\%$ to a substantial $32.2\%$ in an otherwise
identical experiment (\S\ref{sec:utility-results}).

\paragraph{Contributions.} Concretely, this paper:
\begin{enumerate}
\item grounds a scheduling utility function's quality term in
  \emph{measured} downstream semantic-segmentation mIoU rather than pixel
  fidelity, and shows this changes not just the numbers but the
  qualitative conclusion about whether intelligent scheduling matters
  (\S\ref{sec:compression}, \S\ref{sec:utility-results});
\item presents a single \emph{unified utility function} — quality,
  tardiness, resource cost, and a fairness floor, combined so that the
  resulting optimization stays a mixed-integer \emph{linear} program with
  no auxiliary binary variables — that replaces five previously
  disagreeing scoring definitions used across the scheduler, the offline
  bound, and two prior hierarchical variants (\S\ref{sec:formulation});
\item evaluates ten schedulers — four fixed-depth and one adaptive greedy
  baseline, a flat rolling-horizon MPC, a congestion-aware routing
  variant, and two ways of coupling a dedicated routing optimizer to the
  depth decision (hierarchical and two-level/peer) — against a
  three-tiered offline optimality bound, under two link-capacity regimes
  designed to make different resources binding (\S\ref{sec:setup},
  \S\ref{sec:utility-results});
\item shows, from the resulting depth-selection decisions and a
  link-rate sensitivity sweep, that \emph{admission control} — deciding
  early to drop a task rather than let it starve — is what separates the
  strongest and weakest MPC variants, not how cleverly the kept tasks are
  routed (\S\ref{sec:decisions}, \S\ref{sec:sweeps}); and
\item reports a structural, not tunable, negative result: freezing a
  concrete route for a multi-minute planning epoch is usable at its own
  executed slot only a small fraction of the time, because ground-station
  contact windows in this constellation are shorter than any freeze period
  worth having (\S\ref{sec:formulation}, \S\ref{sec:utility-results}); and
  we additionally evaluate a learned (PPO) scheduling policy as a
  compute-cost point of comparison to MPC (\S\ref{sec:rl}).
\end{enumerate}

\paragraph{Roadmap.} \S\ref{sec:related-work} situates this work relative
to orbital edge computing, learned image compression, and MPC-based
network scheduling. \S\ref{sec:formulation} gives the system model and the
full scheduling formulation, including every MPC variant evaluated later.
\S\ref{sec:setup} specifies the simulated constellation, link-capacity
regimes, and task workload used to reproduce every subsequent result.
\S\ref{sec:compression}--\S\ref{sec:rl} present results: compression and
downstream-task quality, the unified-utility scheduler comparison,
depth-selection decision distributions, link-rate sensitivity, and the
learned-policy baseline, in that order. \S\ref{sec:conclusion} summarizes
and discusses what remains open.

\section{System Model and Problem Formulation}
\label{sec:formulation}

This section states the system model and scheduling formulation this project implements. It builds directly on the original system-and-task model and rolling-horizon MPC objective (satellites $\mathcal{S}$, GBSs $\mathcal{G}$, tasks $\mathcal{K}$, depths $\mathcal{D}$, and the basic MPC structure), extending it with a concrete quality signal, a unified multi-term utility, and additional scheduler variants -- congestion-aware routing, two-level and hierarchical MPC couplings, and an offline optimality bound -- that later sections evaluate. Every symbol below matches the implementation (\texttt{hypatia\_sim/oec\_sim/}) exactly; anywhere an earlier design intent and the shipped code differ, it's flagged inline instead of silently resolved.

\subsection{Sets, time, and links}

We consider a LEO orbital edge computing network consisting of a set of satellites $\mathcal{S}$ and a set of ground base stations (GBSs) $\mathcal{G}$. Time is divided into equal scheduling slots $t \in \mathcal{T}$, each of duration $\Delta t$. Because satellites move, the network topology is time-varying: let $\mathcal{E}_t$ denote the set of available inter-satellite links (ISLs) and satellite-to-ground links at slot $t$, where each available link $(i,j) \in \mathcal{E}_t$ has capacity $C_{ij}(t)$, obtainable from predicted contact windows and link states. In this paper, $\mathcal{S}$ is a $1{,}156$-satellite Kuiper-630-shell Walker constellation, $\mathcal{G}$ is a 4-city ground segment, $\Delta t = 30$\,s, and $\mathcal{E}_t$ is gated by a line-of-sight/range test for ISLs and an elevation test for ground links -- full numeric values are in Table~\ref{tab:constellation} (\S\ref{sec:setup}).

\subsection{Tasks and compression depths}

A set of image-transmission tasks $\mathcal{K}$ is generated when satellites observe fixed areas of interest. Each task $k \in \mathcal{K}$ is characterized by its source satellite $s_k$, destination GBS $g_k$, arrival time $a_k$, number of images $N_k$, importance weight $w_k$, and soft deadline $d_k$. Each task is compressed using one effective RQ-VAE depth selected from
\begin{equation}
\mathcal{D} = \{2, 4, 8, 16\}.
\end{equation}
That's the whole choice a task makes: one of four compression strengths, trading payload size against how useful the reconstructed image ends up being.

For each depth $q \in \mathcal{D}$, let $b_q$ denote the compressed payload size per image and $t_q^{\mathrm{enc}}$ the encoding delay. Rather than leaving downstream utility as an abstract placeholder $u_q$, we measure it: $b_q = 88q$ bytes (measured on the FLAIR-1 $8{\times}8$ dedicated models), and quality is either the legacy pixel-fidelity proxy $u_q = 1 - \mathrm{LPIPS}_q$ or the measured downstream segmentation quality $s_q$, both selectable via a configuration flag (which one is preferable, and why, is a results-section question for later):
\begin{equation}
s_q = \frac{\mathrm{mIoU}_q - \mathrm{mIoU}_{\mathrm{floor}}}
           {\mathrm{mIoU}_{\mathrm{ref}} - \mathrm{mIoU}_{\mathrm{floor}}},
\end{equation}
where $\mathrm{mIoU}_{\mathrm{floor}}$ is the downstream segmentation model's score with the optical bands blanked out (near-infrared and elevation channels only) -- the decision-theoretic zero for not delivering an image's optical content, rather than the mathematical floor $\mathrm{mIoU} = 0$. So $s_q = 0$ means "as useless as not sending it," $s_q = 1$ means "as useful as the uncompressed original," and $s_q$ in between reads directly as a fraction of the achievable task value recovered at depth $q$.

Let $x_{k,q} \in \{0,1\}$ indicate whether task $k$ is compressed using depth $q$. Each task selects exactly one compression depth:
\begin{equation}
\sum_{q \in \mathcal{D}} x_{k,q} = 1, \qquad \forall k \in \mathcal{K}.
\end{equation}
Accordingly, the total compressed payload size of task $k$ is
\begin{equation}
B_k = N_k \sum_{q \in \mathcal{D}} b_q\, x_{k,q}.
\end{equation}
Pick one depth per task, and the task's total size is its image count times whichever per-image payload that depth costs -- the accounting every later constraint and objective term builds on.

Let $r_k(t)$ denote the number of images of task $k$ fully delivered to its destination GBS by the end of slot $t$. The delivery fraction of task $k$ is
\begin{equation}
\rho_k(t) = \frac{r_k(t)}{N_k}, \qquad 0 \le \rho_k(t) \le 1.
\end{equation}
Let $y_{k,ij}(t) \ge 0$ denote the number of bits of task $k$ transmitted over link $(i,j)$ during slot $t$. The link-capacity constraint is
\begin{equation}
\sum_{k \in \mathcal{K}} y_{k,ij}(t) \le C_{ij}(t)\, \Delta t,
\qquad \forall (i,j) \in \mathcal{E}_t,
\end{equation}
i.e.\ every task sharing a link has to fit inside that link's capacity for the slot, together.

Let $Q_{i,k}(t)$ denote the remaining amount of task $k$'s data stored at node $i$ at slot $t$. The queue state evolves according to arriving compressed data and transmitted traffic, subject to standard flow conservation over the time-varying network: data leaves a queue only by being transmitted along a feasible link, and can't exceed what has actually arrived at that node so far.

To capture latency sensitivity, let $\phi_k(t) \in [0,1]$ denote the freshness weight of task $k$ at slot $t$:
\begin{equation}
\phi_k(t) =
\begin{cases}
1, & t \Delta t \le d_k, \\
\exp\!\big(-\alpha_k (t \Delta t - d_k)\big), & t \Delta t > d_k,
\end{cases}
\end{equation}
where $\alpha_k$ controls how quickly value decays after the deadline ($\alpha = 1/300\,\mathrm{s}^{-1}$ in our experiments). Delivered before the deadline, an image counts at full value; delivered late, it's still worth something but decays exponentially -- a wildfire photo a day late is close to worthless, one five minutes late is nearly as good as on time. Tasks $|\mathcal{K}| = 64$ arrive from 8 areas of interest with $N_k \sim \mathcal{U}\{100{,}000, 300{,}000\}$ and $w_k \sim \mathcal{U}\{1,2,3\}$ (seed 42); a task may be explicitly \texttt{dropped} (abandoned mid-flight) or \texttt{rejected} (refused at admission), both first-class reported outcomes rather than silent non-delivery.

\subsection{Unified utility}
\label{sec:unified-utility}

Utility used to be defined in five places that disagreed: the reported score, the flat MPC's own objective, the two levels of the hierarchical scheduler, and the offline bound each used a different combination of terms. We use one function for all five. With normalized weight $\bar w_k = w_k / \max_j w_j$ and realized arrival time $t^{\mathrm{arr}} = (t{+}1)\Delta t + \delta_{k,t}$ along the path actually used for delivery,
\begin{align}
U_k &= \bar w_k \Big[\, \omega_Q\, s_{q_k}\, \hat G_k
       \;-\; \omega_T\, T_k \;-\; \omega_E\, C_k \,\Big], \\
U &= \sum_k U_k \;+\; \omega_F\, |\mathcal{K}|\, \min_k \hat U_k,
\end{align}
with default weights $\omega_Q = 1.0$, $\omega_T = 0.25$, $\omega_E = 0.05$, $\omega_F = 0.1$. Each task earns reward proportional to how much of it got delivered, scaled by compression quality ($s_{q_k}$) and delivered fraction ($\hat G_k$), minus a lateness penalty ($T_k$) and a resource-cost penalty ($C_k$); the system-wide score sums every task's net reward and adds a bonus tied to whichever task did worst, so nothing can rack up a high score by quietly abandoning one inconvenient task. Here $\hat G_k \in [0,1]$ is a concave coverage gain over delivered, freshness-weighted image fraction (\S\ref{sec:coverage}); tardiness is
\begin{equation}
T_k = \sum_t \frac{\Delta r_{k,t}}{N_k}
  \min\!\Big(1,\, \frac{\max(0,\, t^{\mathrm{arr}} - d_k)}{\Delta_{\mathrm{ref}}}\Big),
\end{equation}
so each newly-delivered slice of images is penalized in proportion to how late it arrived past the deadline, capped at 1; resource cost is
\begin{equation}
C_k = \sum_t \frac{\Delta r_{k,t}}{N_k}
  \Big[\hat\omega_{\mathrm{tx}} \frac{b_q}{b_{\max}} + \hat\omega_{\mathrm{enc}}\Big],
\end{equation}
a normalized charge for bandwidth and encode effort spent on each delivered slice; and $\hat U_k = U_k / (\omega_Q\, s_{\max}\, \bar w_k) \le 1$ is a scale-free, weight-relative normalization of $U_k$ used only inside the fairness term, so a naturally small task (low weight) doesn't get unfairly flagged as the worst-off one just because its raw $U_k$ is small.

\textbf{Legacy identity.} Set $\omega_Q{=}1$, $\omega_T{=}\omega_E{=}\omega_F{=}0$, one coverage segment of unit width/slope, no weight normalization, and $s_q = 1 - \mathrm{LPIPS}_q$, and this collapses exactly to the original score $\sum_k \sum_t w_k \phi_k(t^{\mathrm{arr}}) u_{q_k} \Delta r_{k,t} / N_k$ -- legacy scoring is a parameter setting of this same function, not a separate code path, and every previously-committed legacy number reproduces exactly under this setting.

\subsection{Concave coverage without integer variables}
\label{sec:coverage}

Getting the first fraction of a task's images delivered is worth more, per image, than getting the last fraction -- a concave, diminishing-returns reward. Encoding a curved reward inside a linear program normally needs extra machinery (binary or SOS2 variables, expensive to solve). Instead, coverage gain is linearized over $J{=}4$ segments of width $\Delta_j$ and strictly decreasing slopes $m_1 > \cdots > m_J$ with $\sum_j \Delta_j m_j = 1$ (default: widths $(0.25,0.25,0.25,0.25)$, slopes $(2.00, 1.12, 0.60, 0.28)$, a discretization of $g(\rho) = (1-e^{-3\rho})/(1-e^{-3})$):
\begin{align}
\hat G_k &= \sum_{\tau, j} m_j\, \phi_k(t^{\mathrm{arr}}_\tau)\,
  \lambda_{k,q,\tau,j}, \\
\sum_j \lambda_{k,q,\tau,j} &= \frac{y_{k,q,\tau}}{N_k}, \qquad
\sum_{q,\tau} \lambda_{k,q,\tau,j} \le \Delta_j^{\mathrm{res}},
\end{align}
where $\Delta_j^{\mathrm{res}}$ is the width of segment $j$ still unclaimed by coverage already realized before the current planning horizon (skip this residual and the MPC would re-earn steep early-segment credit at every re-plan, double-counting value already banked). Picture each task's 0--100\% delivered range as four stacked buckets, each paying less per percentage point than the one before; $\lambda_{k,q,\tau,j}$ is how much of bucket $j$ a given slot's delivery fills. Because the payoff always shrinks bucket by bucket, a solver that only understands straight lines will still fill bucket 1 before touching bucket 2, entirely on its own -- no binary "which bucket am I in" variable needed, and the linear relaxation of this objective is exact.

\subsection{Fairness: maximin, not Jain}

Jain's index $(\sum_k U_k)^2 / (|\mathcal{K}| \sum_k U_k^2)$ is a classic networking fairness statistic, but it's a ratio of quadratics and can't be represented inside a mixed-integer linear program, so it's reported only as a diagnostic, never optimized directly. What's actually optimized is a maximin floor -- how well the single worst-off task did -- via one continuous column $u_{\min}$ and $|\mathcal{K}|$ constraint rows $u_{\min} \le \hat U_k$, entering the objective as $\omega_F |\mathcal{K}| u_{\min}$ (the second term of $U$ above). $u_{\min}$ needs a free, unbounded-below lower bound, since $\hat U_k$ can go negative once tardiness/cost penalties exceed the quality gain for a badly-served task.

\subsection{Flat MPC objective}
\label{sec:flat-mpc}

At each decision slot $t$, the scheduler observes the current network state -- task queues, remaining deadlines, satellite positions, available links -- and predicts network evolution over a finite horizon of $H$ slots: $\hat{\mathcal{E}}_{\tau|t}$, $\hat C_{ij}(\tau|t)$, $\hat Q_{i,k}(\tau|t)$ for $\tau = t, \ldots, t+H-1$. It jointly determines compression depth, transmission timing, and traffic allocation over that horizon, then executes only the first-slot decisions; at slot $t+1$ it observes the new state, shifts the horizon forward, and solves again. The rolling-horizon MILP solved by the flat \texttt{mpc} scheduler (re-solved every 5 slots or on new-task arrival; $H = 60$ slots) is
\begin{equation}
\max \sum_{k,q,\tau} \Big[\, w_k\, \phi_k(t^{\mathrm{arr}}_{k,\tau})\, u_q
  \;+\; \lambda_1 (H-\tau)\, b_q \cdot 8 / 10^9
  \;-\; \lambda_{\mathrm{late}}\, w_k\,
    \frac{\max(0,\, t^{\mathrm{arr}}_{k,\tau} - d_k)}{\Delta_{\mathrm{ref}}}
  \,\Big] \frac{y_{k,q,\tau}}{N_k},
\end{equation}
with $\lambda_{\mathrm{late}} = 5\times10^{-2}$, $\Delta_{\mathrm{ref}} = 300$\,s, subject to: one depth per task per slot, link/GBS capacity, the encoder throughput pipeline (81 images/s per satellite), and cumulative delivered-volume consistency. Here $t^{\mathrm{arr}}_{k,\tau} = (t+\tau+1)\Delta t + \text{route delay}$ uses the real propagation delay of the path actually assigned to $(k,\tau)$. The first bracket term rewards weighted, freshness-discounted useful delivery; the second gives a small forward-looking bonus to plans that keep transmitting rather than idle capacity; the third subtracts an explicit lateness penalty -- which lets the MPC rationally choose to stop pushing a task's backlog once its cost outweighs its already-decayed value, instead of that decision happening implicitly through silent non-delivery. \S\ref{sec:hier} below makes this same tradeoff explicit via admission control and drop.

\subsection{Congestion-aware predictive routing}
\label{sec:congestion-routing}

The default router picks the geometrically shortest path regardless of how busy it currently is, which can overload a popular link while a nearly-as-fast alternative sits idle. The congestion-aware variant (\texttt{mpc-congestion}) instead predicts each link's cost as propagation delay plus a queueing-theory-style congestion penalty computed against the MPC's own current plan:
\begin{equation}
w_e(\tau) = d_e(\tau) + \beta\, d_{\mathrm{ref}}\,
  \frac{\rho_e(\tau)}{1 - \min(\rho_e(\tau),\, \rho_{\max})}, \qquad
\rho_e(\tau) = \frac{f_e(\tau)}{C_e \Delta t},
\end{equation}
where $f_e(\tau)$ is the bits the current plan places on edge $e$ at predicted step $\tau$, and $\rho_e(\tau)$ is that edge's predicted utilization. A link's effective cost stays close to its physical delay while mostly empty, then rises sharply as it approaches saturation ($\rho_e \to \rho_{\max}$), which spreads traffic across several good paths instead of jamming everyone onto the one "best on paper" route. Because the plan depends on routes and routes depend on the plan, this is iterated to a damped fixed point (method-of-successive-averages, $\eta_{\mathrm{it}} = 1/(\mathrm{it}+2)$, 3 rounds, keep-best-iterate), with $\rho$ clipped at 0.99 so a saturated link is never treated as disconnected. This coupling is only measurable under the fabric-limited capacity regime (\S\ref{sec:setup}) -- under the gbs-limited regime the ground-station aggregate budget dominates every route regardless of ISL cost, and congestion-aware routing is a provable, verified no-op there.

\subsection{Two-MPC couplings: routing and depth}

\textbf{\texttt{mpc-2level}}: a dedicated routing MPC solves a min-cost multicommodity-flow linear program over candidate paths,
\begin{equation}
\min \sum_{k,p,\tau} f_{k,p,\tau}
  \frac{\delta_{p,\tau}}{\Delta^{\mathrm{delay}}_{\mathrm{ref}}}
  + M \sum_{k,\tau} \mathrm{sh}_{k,\tau}
\quad \text{s.t.} \quad
\sum_p f_{k,p,\tau} + \mathrm{sh}_{k,\tau} = D_{k,\tau}, \quad
\sum_k \sum_{p \ni e} f_{k,p,\tau} \le C_e \Delta t,
\end{equation}
where $f_{k,p,\tau}$ is the flow assigned to task $k$ on candidate path $p$, $\mathrm{sh}_{k,\tau}$ is an explicit "shortfall" variable absorbing any demand no path can carry (penalized heavily by $M$ so it's only used when genuinely unavoidable), and $D_{k,\tau}$ is task $k$'s desired bit demand at step $\tau$. This LP spreads each task's traffic across whichever candidate paths have room, at minimum total delay, instead of committing everything to one path -- and if no combination of paths can carry all of it, the shortfall variable makes that gap explicit rather than the LP going infeasible. This routing LP and the depth MILP (\S\ref{sec:flat-mpc}) exchange demand $D_{k,\tau}$ and per-edge flow shares, iterated to a fixed point with the same successive-averages damping as \S\ref{sec:congestion-routing}. Critically, $D_{k,\tau}$ has to be the task's \emph{desired} bits, not what the depth MILP already conceded -- deriving it from the previous plan would be circular and silently reproduce flat MPC even under real congestion (a correctness trap caught during development). Iteration 0 is defined to be exactly flat \texttt{mpc}, so \texttt{mpc-2level} can never plan worse than flat MPC plans, by construction.

\textbf{\texttt{mpc-hier-route}}: a route-coordinator MPC runs once per 10-minute macro-epoch over a 60-slot horizon in 10-slot macro-windows (topology evaluated at each window's middle slot, as a forecast), solving the same flow LP to freeze the top-$K$ candidate paths per task for the whole epoch; the fast depth MILP then solves only within those frozen paths. Empty route assignments reproduce \texttt{mpc-hier} (\S\ref{sec:hier}) exactly, verified bit-for-bit.

\textbf{Measured: route freezing does not survive a LEO fabric.} A route frozen at epoch start is usable at its own executed slot only $14.4\%$ of the time in general (our fabric-limited scenario: $7.7\%$), falling to $0\%$ by a 570-second lookahead -- because ground-station contact windows average only 227--265\,s, shorter than any planning epoch worth having. That's a structural property of the constellation geometry, not a tuning failure: it's a negative result about freezing \emph{concrete paths} specifically, not necessarily about freezing a more abstract routing decision (e.g.\ a serving-GBS assignment), which remains untested.

\subsection{Hierarchical MPC and the rotting-backlog fix}
\label{sec:hier}

\texttt{mpc-hier} pairs a slow upper-level coordinator (cadence 20 slots / 10\,min, horizon 6 macro-steps $\approx$ 2\,h, an LP over continuous relaxed depth $x_{k,q} \in [0,1]$ and GBS-aggregate budget only) with a fast lower-level per-task MILP (cadence 5 slots or on arrival, horizon 20 slots, integral depth $x_{k,q}$, the same machinery as \S\ref{sec:flat-mpc}). The upper level decides, at a coarse cadence, which tasks get admitted and roughly how much budget each gets; the lower level decides, at a fine cadence, the exact depth and schedule within that budget -- a smaller, faster-solved problem in exchange for real admission/drop activity. Three mechanisms address a specific failure mode, stale-backlog starvation (a task whose freshness has decayed to near zero but still generates optimization variables and is never served nor removed): admission control at arrival, explicit drop once freshness falls below a threshold or a task is sufficiently overdue (reported as \texttt{dropped}, not hidden), and a backpressure aging term in the lower-level objective that raises priority the longer a task has waited. These thresholds are first-pass tuned (five combinations tried), not an exhaustive search.

\textbf{\texttt{mpc-hier} vs. the routing/depth split.} \texttt{mpc-hier}'s two levels both operate on depth (a relaxed depth budget above, an integral depth choice below), with static routing at both levels -- a different pairing from \texttt{mpc-2level}/\texttt{mpc-hier-route}, whose split is one MPC dedicated to routing and one to depth. \todoconfirm{Xuanhao's email asked to drop "the previous mpc-hier results where the two MPCs are both used for depth selection." That description textually matches \texttt{mpc-hier} as just defined, which makes it the likely candidate for what Xuanhao means -- but this is still pending his direct confirmation of scope.}

\subsection{Offline optimality bound}

To know whether a scheduler's score is actually good, it has to be compared to the best \emph{possible} score under the same rules -- an offline bound, computed three ways of increasing tightness: (1) an analytic ceiling (every image at the best depth, delivered instantly -- a sanity check only, not a meaningful bound); (2) a linear-programming bound (depth relaxed to continuous $x_{k,q} \in [0,1]$, time aggregated into 5-slot windows, provably-non-binding link/GBS rows dropped, utility evaluated at each window's \emph{start} to remain a genuine upper bound -- never optimistic about when delivery actually happens); and (3) a mixed-integer bound (the same relaxed constraints at full time resolution with integral depth, reporting the solver's dual bound even when the search hasn't fully converged within its time budget, which remains a valid upper bound regardless). The governing correctness rule: every term in a scheduler's realized score must also appear in the bound's objective, relaxed only in the optimistic direction -- enforced at runtime by an assertion that the bound is never smaller than any scheduler's realized score.

\section{Experimental Setup}
\label{sec:setup}

This section specifies every parameter needed to reproduce the results in
this paper from the same simulator (\texttt{hypatia\_sim/oec\_sim/}, built on
the Hypatia LEO orbital-mechanics/topology engine).

\subsection{Constellation and ground segment}

Table~\ref{tab:constellation} lists the physical scenario parameters
(\texttt{hypatia\_sim/oec\_sim/config.py}). Four named constellations are
implemented (\texttt{kuiper-630}, \texttt{starlink-550}, \texttt{telesat-1015},
\texttt{small-walker}); all results in this paper use \texttt{kuiper-630}
unless stated otherwise.

\begin{table}[h]
\centering
\caption{Constellation, link, and task-load parameters.}
\label{tab:constellation}
\begin{tabular}{ll}
\toprule
Parameter & Value \\
\midrule
Constellation & Kuiper-630 shell 1: Walker-delta $1156/34/1$ (34 planes $\times$ 34 sats) \\
Altitude / inclination & 630\,km / $51.9^\circ$ \\
Orbital period & 5{,}830\,s \\
Ground stations $G$ & Tokyo, New York, S\~ao Paulo, Sydney (4 total), min.\ elevation $20^\circ$ \\
ISL topology & +Grid, 2{,}312 candidate links; feasible if graze altitude $>80$\,km \\
             & above Earth's surface \emph{and} range $\le 5{,}016$\,km \\
Areas of interest & 8 (California-fires, Amazon-basin, Sahel, Ganges-delta, \\
                   & Barrier-reef, Dnipro-basin, and 2 more; see \texttt{config.py:AOIS}) \\
Tasks per run $|K|$ & 64, one per (AOI, publication event), Poisson-triggered \\
Images per task $N_k$ & $\sim \mathcal{U}(100{,}000,\ 300{,}000)$, total $\approx 12.81$M images \\
Task weight $w_k$ & $\in \{1, 2, 3\}$ \\
Random seed & 42 (task generation; sweeps in \S\ref{sec:sweeps} use 5 seeds) \\
Simulation horizon & 18{,}000\,s (5\,h, $\approx 3.1$ orbits) \\
Slot length $\Delta t$ & 30\,s $\Rightarrow$ 601 decision slots \\
MPC planning horizon $H$ & 60 slots (30\,min look-ahead) \\
\bottomrule
\end{tabular}
\end{table}

\subsection{Link capacity regimes}

Two named capacity regimes share the identical topology and task model but
rebalance link rates to make different resources binding
(Table~\ref{tab:regimes}). \textbf{gbs-limited} is the default and is used
for the legacy-vs-unified utility comparison in \S\ref{sec:utility-results};
its per-GBS aggregate budget is $\ge 12\times$ tighter than any ISL segment a
route could cross, so ISL contention is mathematically unreachable there.
\textbf{fabric-limited} throttles the ISL rate specifically so routing
decisions have a measurable effect, and is used for the ten-scheduler
comparison in \S\ref{sec:utility-results} and the decision-distribution
analysis in \S\ref{sec:decisions}.

\begin{table}[h]
\centering
\caption{Link capacity regimes (\texttt{config.py:\_SCENARIOS}).}
\label{tab:regimes}
\begin{tabular}{lrrr}
\toprule
Regime & ISL rate & GSL rate (per sat--GBS pair) & GBS aggregate rate \\
\midrule
gbs-limited (default) & 100\,Mbps & 100\,Mbps & 2\,Mbps \\
fabric-limited         & 1\,Mbps   & 50\,Mbps  & 500\,Mbps \\
\bottomrule
\end{tabular}
\end{table}

\subsection{Compression depths and payload}

Each task-image is encoded on board at one of $D = \{2, 4, 8, 16\}$
RQ-VAE residual-quantization depths before downlink. Payload scales linearly
with depth, $b_q = 88q$ bytes (measured on the FLAIR-1 $8{\times}8$ dedicated
models, \S\ref{sec:compression}): $b_2 = 176$\,B, $b_4 = 352$\,B,
$b_8 = 704$\,B, $b_{16} = 1{,}408$\,B. On-board encode time is
$12.34\,\text{ms}\,\pm\,0.03$--$0.05\,\text{ms}$ per image and is
\emph{depth-independent} (measured identically at $8{\times}8{\times}1$ and
$8{\times}8{\times}8$), giving an encoder throughput of 81 images/s per
satellite regardless of the depth chosen.

\subsection{Decision cadence: what ``every 30\,s'' actually means}

The scheduling slot is $\Delta t = 30$\,s, but the MPC does \emph{not}
re-solve its optimization every slot. Each solve plans $H=60$ slots (30\,min)
ahead and commits only to the first slot of that plan; the scheduler
re-solves only when either 5 slots have elapsed since the last solve, or a
new task arrives — otherwise it executes the next committed step of the
existing plan. Over the full 601-slot (5-hour) simulation, the flat
\texttt{mpc} scheduler solves the optimization approximately 121 times, not
601 times. \textbf{Reported solve-time statistics (mean/p95/max) are for a
single individual solve} (how long the optimizer takes to think once); a
separately reported \emph{total} solve-time figure is the sum of individual
solve times across the whole run — total compute time spent solving, not
wall-clock simulation time and not a per-slot cost.

\subsection{Software and solver}

The scheduling MILPs/LPs (flat MPC, hierarchical admission/budget LP, the
two-level routing LP, and the offline oracle bound) are solved with HiGHS via
\texttt{scipy.optimize.linprog}/\texttt{milp}. Orbital mechanics, contact
windows, and link feasibility come from Hypatia's satellite-generation and
topology modules (\texttt{hypatia/satgenpy}), vendored unmodified.

\subsection{Quality-table provenance}

The utility function's quality term $s_q$ (\S\ref{sec:utility-results},
\S\ref{sec:formulation}) is grounded in measured downstream FLAIR-1
segmentation mIoU, source label \texttt{flair-unet-r34-rgbie / val7050 /
2026-08-31} (full 7{,}050-image validation population; see
\S\ref{sec:compression} for the measurement itself). All results in
\S\ref{sec:utility-results} and \S\ref{sec:decisions} use this real,
measured table (\texttt{hypatia\_sim/oec\_sim/quality\_table.json}), not the
placeholder fallback (\texttt{config.QUALITY\_TABLE\_FALLBACK}) that the code
falls back to when this file is absent.


% Saved for later, not yet included in this draft (files exist in
% sections/, ready to add back once results/related work are finished):
%   \section{Related Work}
\label{sec:related-work}

\todoconfirm{This section is drafted against general background rather
than the two reference papers Xuanhao sent for how experimental results
should be presented — those weren't available in the session that drafted
this paper. Once shared, this section's structure/tone and citation set
should be revised to match them; the citations below are limited to works
this project directly builds on, to avoid citing anything not actually
verified.}

\paragraph{Orbital edge computing and LEO satellite networking.} Moving
computation to the network edge is a well-established idea in terrestrial
systems; OEC applies the same principle where the "edge" is the satellite
itself, motivated by a downlink budget that is orders of magnitude smaller
than the volume of imagery a modern sensor can capture in a single pass.
LEO constellations add a networking wrinkle largely absent from
terrestrial edge computing: topology and link capacity are not merely
variable but \emph{predictable} — satellite and ground-station positions,
and therefore contact windows and candidate routes, are known in advance
from orbital mechanics. This paper's simulation environment is built on
Hypatia~\cite{kassing2020hypatia}, an open-source LEO network simulator
that models constellation geometry, inter-satellite links, and
ground-station contact windows explicitly, which is what makes a
rolling-horizon scheduler that \emph{predicts} short-term future
connectivity (rather than only reacting to it) both possible and testable
at scale.

\paragraph{Learned image compression for downlink-constrained systems.}
On-board compression here uses a residual-quantized variational
autoencoder (RQ-VAE)~\cite{lee2022autoregressive}: an encoder maps an
image to a small set of codebook indices, and successive
\emph{residual} quantization stages progressively refine the
reconstruction, each stage correcting the error left by the stage before
it. The number of stages run before transmission — this paper's
\emph{depth} — is exactly the fidelity/payload knob a scheduler has to
reason about. We evaluate this compression model on two remote-sensing
benchmarks: EuroSAT~\cite{helber2019eurosat}, a 10-class land-cover
classification dataset used for an initial spatial-size/depth sweep, and
FLAIR-1~\cite{garioud2022flair}, a larger aerial-imagery semantic-segmentation
dataset (15 land-cover classes) used for the paper's central downstream-task
evaluation. Unlike most learned-compression evaluations, which report
pixel-fidelity metrics (PSNR, SSIM, LPIPS, FID) as an end in themselves,
this paper treats those metrics as necessary but insufficient: our results
(\S\ref{sec:downstream}) show they can be nearly flat exactly where a
downstream task's performance is not.

\paragraph{MPC-based scheduling and network optimization.} Model
Predictive Control — solve an optimization over a finite look-ahead
horizon, act on only the first step, then re-plan as new information
arrives — is a natural fit for a system whose future state (contact
windows, link capacities) is partially predictable but subject to real
uncertainty (bursty task arrival, imperfect topology forecasts). We
formulate the scheduling problem as a mixed-integer linear program so it
remains tractable at the problem sizes this paper evaluates, and construct
a piecewise-linear coverage reward specifically so that a concave,
diminishing-returns objective can be optimized without introducing
additional binary or SOS2 variables (\S\ref{sec:coverage}) — a standard
technique for keeping a naturally nonlinear objective inside a linear
program's reach. We evaluate this rolling-horizon approach against both
simpler greedy baselines and a learned (reinforcement-learning) policy
(\S\ref{sec:rl}), and against hierarchical and peer couplings that split
routing and depth decisions into separate, communicating optimizers
(\S\ref{sec:formulation}).

%   \section{Compression Model Results on FLAIR-1}
\label{sec:compression}

This section reports the on-board RQ-VAE compression model's reconstruction
quality, the downstream segmentation task performance it induces (which
grounds the utility function's quality term $s_q$), and on-board compute
profiling.

\subsection{Reconstruction quality: dedicated models}

Table~\ref{tab:flair-dedicated} reports reconstruction quality for
dedicated RQ-VAE models trained at each depth, evaluated on the full
7{,}050-image FLAIR-1 validation split, fixed $8\times 8$ spatial latent grid.
Depths 2, 3, and 4 have no dedicated trained model (training configs exist
but were not run); \S\ref{sec:truncation} covers how those conditions are
obtained instead.

\begin{table}[h]
\centering
\caption{Dedicated FLAIR-1 RQ-VAE models, full validation split.}
\label{tab:flair-dedicated}
\begin{tabular}{lrrrrrr}
\toprule
Model & Payload & PSNR (dB) & SSIM & LPIPS & FID & Compression \\
\midrule
$8{\times}8{\times}1$  & 88\,B    & 20.63 & 0.4560 & 0.4595 & 71.33  & 8{,}937:1 \\
$8{\times}8{\times}8$  & 704\,B   & 21.02 & 0.4643 & 0.4779 & 73.19  & 1{,}117:1 \\
$8{\times}8{\times}16$ & 1{,}408\,B & 21.06 & 0.4899 & 0.5039 & 125.78 & 559:1 \\
\bottomrule
\end{tabular}
\end{table}

\subsection{The truncation experiment: is reuse viable?}
\label{sec:truncation}

To avoid training a dedicated model per depth, we tested whether a single
depth-16 model, decoded using only its first $N$ of 16 residual stages at
inference (\texttt{forward\_partial\_code}), can substitute for a dedicated
depth-$N$ model. Table~\ref{tab:truncation} reports this truncated-inference
condition against the same depth-16 checkpoint.

\begin{table}[h]
\centering
\caption{Truncated inference from one depth-16 checkpoint.}
\label{tab:truncation}
\begin{tabular}{lrrrr}
\toprule
Truncated to & PSNR (dB) & SSIM & LPIPS & FID \\
\midrule
1 stage   & 18.67 & 0.4793 & 0.5948 & 225.25 \\
2 stages  & 19.72 & 0.4838 & 0.5570 & 197.76 \\
4 stages  & 20.35 & 0.4831 & 0.5314 & 163.90 \\
8 stages  & 20.73 & 0.4861 & 0.5171 & 143.19 \\
16 stages & 21.06 & 0.4899 & 0.5039 & 125.78 \\
\bottomrule
\end{tabular}
\end{table}

\textbf{Result: reuse is not viable.} At depth 1, truncated inference is
2\,dB worse in PSNR and over $3\times$ worse in FID than the dedicated
depth-1 model (225.25 vs.\ 71.33). Even at depth 8, where the gap narrows,
truncated inference remains measurably behind dedicated training (FID 143.19
vs.\ 73.19). Consequently, the depth-2/4/8 downstream conditions reported in
\S\ref{sec:downstream} use truncated inference from the depth-16 checkpoint
and should be read as a \emph{lower bound} on what a dedicated model at that
depth would achieve, not its true ceiling.

\subsection{Downstream task performance}
\label{sec:downstream}

Pixel-fidelity metrics (PSNR, LPIPS) do not measure whether a reconstruction
is still useful for the mission task. We evaluate each compression depth's
reconstructions on FLAIR-1 semantic segmentation (15 land-cover classes)
using IGNF's pretrained \texttt{FLAIR-INC\_rgbie\_15cl\_resnet34-unet} model,
scored by mean IoU (mIoU) over the full 7{,}050-image validation population
(Table~\ref{tab:downstream}).

\begin{table}[h]
\centering
\caption{Downstream segmentation mIoU by compression depth. $s_q$ is
floor-anchored: $s_q = (\text{mIoU}_q - \text{mIoU}_{\text{floor}}) /
(\text{mIoU}_{\text{ref}} - \text{mIoU}_{\text{floor}})$.}
\label{tab:downstream}
\begin{tabular}{lrr}
\toprule
Condition & mIoU & $s_q$ (floor-anchored) \\
\midrule
blank (NIR + elevation only, floor) & 6.08\% & --- \\
$q=1$ (truncated)  & 12.60\% & 0.104 \\
$q=2$ (truncated)  & 17.53\% & 0.182 \\
$q=4$ (truncated)  & 22.54\% & 0.262 \\
$q=8$ (truncated)  & 25.68\% & 0.312 \\
$q=16$ (dedicated) & 28.70\% & 0.360 \\
uncompressed (ref) & 68.87\% & --- \\
\bottomrule
\end{tabular}
\end{table}

This measured spread ($s_q$: $+247\%$ from $q=1$ to $q=16$) is far wider
than the pixel-fidelity proxy it replaces ($1-\text{LPIPS}$: only $+12\%$
over the same range), and is the reason the unified utility function's
quality term is grounded in mIoU rather than pixel fidelity
(\S\ref{sec:formulation}).

Per-class results show this recovery is uneven: fine-texture classes
(swimming pools, coniferous/deciduous forest) lose 70--99\% of their
achievable IoU at $q=1$, while broad classes (agricultural land,
greenhouses) lose only 28--35\%. Per-class detail is reported as
supplementary analysis; the utility function uses overall mIoU by design.

\todoconfirm{The uncompressed-reference mIoU (68.87\%) is higher than both
the checkpoint's expected reproduction target ($\approx 54.43\%$) and its
Hugging Face model-card self-reported value (58.6\%). A domain-leakage check
between the validation population and the training domain came back clean,
and per-class scores are structurally sane, but this has not been verified
against FLAIR-1's true held-out test split. Treat 68.87\% as internally
consistent for comparing depths against each other, not as a verified match
to an external published number.}

\subsection{On-board compute profiling}

Encode time was measured on an NVIDIA A6000 at $8{\times}8{\times}1$ and
$8{\times}8{\times}8$: $12.34\,\text{ms} \pm 0.03$--$0.05\,\text{ms}$ per
image at both depths, i.e.\ encode time is depth-independent (dominated by
the encoder CNN forward pass, not the residual-quantization loop). Per-image
energy at 30\,W (Jetson AGX Orin mid-band assumption, since the A6000 used
for timing is not a flight part): encode $= 0.370$\,J; downlink at $q=16$,
$11{,}264$ bits at $2\times10^{-7}$ J/bit $= 2.25\times10^{-3}$\,J. Encode
dominates downlink energy by $\sim 164\times$ per image ($\sim 230\times$
over a run's full depth mix), and since it is depth-independent, a
Joule-denominated cost term cannot drive depth choice — depth selection in
this system is a \emph{bandwidth} decision, not an energy one
(\S\ref{sec:formulation} discusses how the utility's cost term is scaled
accordingly).

%   \section{Unified Utility: Baselines, Method, and Variants}
\label{sec:utility-results}

\subsection{Legacy vs.\ unified utility: not the same scale}

Two utility functions are implemented, selected by \texttt{--utility
legacy|unified}. \textbf{Legacy} is a single weighted quality-times-freshness
sum with no tardiness, cost, or fairness term, using a pixel-fidelity
quality proxy ($1-\text{LPIPS}$). \textbf{Unified} adds a tardiness penalty,
a resource-cost penalty, and a maximin fairness floor, and replaces the
proxy with the measured downstream mIoU term $s_q$ from
\S\ref{sec:downstream} (full formulation in \S\ref{sec:formulation}). These
two modes are \emph{not} on the same scale and must not be compared directly
in a single table; Table~\ref{tab:legacy-vs-unified} reports each mode's
own best-fixed-depth-greedy vs.\ MPC margin, from the \texttt{gbs-limited}
scenario, to show how switching the quality term changes the conclusion.

\begin{table}[h]
\centering
\caption{MPC vs.\ best fixed-depth greedy, gbs-limited scenario, seed 42.}
\label{tab:legacy-vs-unified}
\begin{tabular}{lrrr}
\toprule
Utility mode & \texttt{mpc} & best fixed greedy & MPC margin \\
\midrule
legacy ($1-\text{LPIPS}$ quality) & 64.95 & 64.15 (fixed-8) & $+1.2\%$ \\
unified (real mIoU quality)       & 18.13 & 13.71 (fixed-8) & $+32.2\%$ \\
\bottomrule
\end{tabular}
\end{table}

Under the pixel-fidelity proxy, MPC essentially ties the best fixed-depth
greedy baseline. Once quality reflects the real downstream task, MPC's
percentage margin is $\approx 26\times$ larger. This is the central result
motivating the unified utility: the legacy proxy under-states MPC's value
because it is nearly flat across compression depths (\S\ref{sec:downstream}).

\subsection{Compared algorithms}

\begin{itemize}
\item \textbf{greedy-fixed-\{2,4,8,16\}}: always encode at a fixed depth,
  transmit whatever fits in the current slot's remaining capacity.
\item \textbf{greedy-adaptive}: greedy, but adapts encode depth to queue
  length.
\item \textbf{mpc}: flat rolling-horizon MILP (\S\ref{sec:formulation}),
  static shortest-path routing.
\item \textbf{mpc-congestion}: flat MPC with predictive, congestion-aware
  routing (M/M/1-style edge cost, iterated to a fixed point).
\item \textbf{mpc-hier}\todoconfirm{scope of this variant to keep/drop is
  pending Xuanhao's confirmation, see note below}: two-level hierarchy —
  a slow admission/budget LP plus a fast per-task depth MILP, with explicit
  task drop/abandon logic.
\item \textbf{mpc-hier-route}: \texttt{mpc-hier} plus a routing MPC that
  freezes a concrete path per 10-minute macro-epoch.
\item \textbf{mpc-2level}: a routing MPC (multicommodity-flow LP over
  candidate paths) and the depth MPC iterated to a fixed point as
  \emph{peers}, not hierarchically.
\end{itemize}

\todoconfirm{Xuanhao's email asked to drop ``the previous mpc-hier results
where the two MPCs are both used for depth selection.'' \texttt{FORMULATION.md}
describes \texttt{mpc-hier}'s two levels as: upper = LP over a
\emph{continuous relaxed depth} $x_{k,q}\in[0,1]$ (for budgeting), lower =
MILP over \emph{integral depth} $x_{k,q}$ -- i.e.\ both levels of
\texttt{mpc-hier} specifically operate on depth, with static routing at
both. This is different from \texttt{mpc-hier-route}/\texttt{mpc-2level},
which are an explicit routing-MPC + depth-MPC split. This textual match
makes \texttt{mpc-hier} the likely candidate for what Xuanhao means, but all
schedulers are reported below pending his direct confirmation --
see \S\ref{sec:hier} for the full two-level description.}

\subsection{Ten schedulers, one run}

Table~\ref{tab:scheduler-comparison} reports all ten schedulers on the
\texttt{fabric-limited} scenario (where routing choices are measurable),
unified utility, real mIoU quality table, single run/seed (42). Every row is
from the same run, so utility, delivery, on-time rate, and bound-gap are
directly comparable across schedulers.

\begin{table}[h]
\centering
\caption{Scheduler comparison, fabric-limited scenario, seed 42. Bound
(offline HiGHS upper bound, \S\ref{sec:formulation}) $= 18.99$.}
\label{tab:scheduler-comparison}
\begin{tabular}{lrrrr}
\toprule
Scheduler & Utility & Delivery\% & On-time\% & Gap to bound \\
\midrule
greedy-fixed-2   & 6.43  & 100\% & 88.7\% & 66.1\% \\
greedy-fixed-4   & 11.08 & 100\% & 88.7\% & 41.7\% \\
greedy-fixed-8   & 13.71 & 100\% & 88.6\% & 27.8\% \\
greedy-fixed-16  & 8.80  & 100\% & 81.5\% & 53.7\% \\
greedy-adaptive  & 8.80  & 100\% & 81.5\% & 53.7\% \\
mpc              & 18.76 & 87.9\% & 86.3\% & 1.2\% \\
mpc-congestion   & 18.86 & 89.3\% & 87.6\% & 0.7\% \\
mpc-hier         & 14.76 & 78.7\% & 75.0\% (39 dropped) & 22.3\% \\
mpc-hier-route   & 14.78 & 81.9\% & 78.6\% (40 dropped) & 22.2\% \\
\textbf{mpc-2level} & \textbf{18.90} & \textbf{90.3\%} & \textbf{88.6\%} & \textbf{0.5\%} \\
\bottomrule
\end{tabular}
\end{table}

\emph{Note:} \texttt{greedy-adaptive}'s row is numerically identical to
\texttt{greedy-fixed-16} in this run — in this scenario, the adaptive
policy's depth choice converges to always selecting depth 16, so this is a
genuine coincidence in the underlying decisions (verified in
\S\ref{sec:decisions}), not a duplicated row.

\subsection{Individual utility components}

Table~\ref{tab:utility-components} decomposes total utility into its four
additive terms, from the same run as Table~\ref{tab:scheduler-comparison}
(\texttt{hypatia\_sim/oec\_scenario\_couplings/summary.txt}).
$U = \text{quality} - \text{tardy} - \text{cost} + \text{fair}$; $u_{\min}$
is the optimized maximin fairness floor and Jain's index is reported as a
diagnostic only (not MILP-representable, so not itself optimized; see
\S\ref{sec:formulation}).

\begin{table}[h]
\centering
\caption{Utility component decomposition, fabric-limited scenario, seed 42.}
\label{tab:utility-components}
\begin{tabular}{lrrrrrrrr}
\toprule
Scheduler & Total & Quality & $-$Tardy & $-$Cost & $+$Fair & $u_{\min}$ & Jain & Enc.\ (kJ) \\
\midrule
mpc              & 18.757 & 16.344 & $-$0.029 & $-$2.121 & 4.564  & 0.7131  & 0.9978 & 4170.2 \\
greedy-adaptive  & 8.797  & 16.170 & $-$1.502 & $-$2.350 & $-$3.521 & $-$0.5502 & 0.9118 & 4742.3 \\
greedy-fixed-2   & 6.431  & 8.414  & $-$0.813 & $-$1.322 & 0.152  & 0.0237  & 0.9331 & 4742.3 \\
greedy-fixed-4   & 11.076 & 12.089 & $-$0.813 & $-$1.469 & 1.269  & 0.1983  & 0.9670 & 4742.3 \\
greedy-fixed-8   & 13.711 & 14.394 & $-$0.814 & $-$1.763 & 1.894  & 0.2959  & 0.9748 & 4742.3 \\
greedy-fixed-16  & 8.797  & 16.170 & $-$1.502 & $-$2.350 & $-$3.521 & $-$0.5502 & 0.9118 & 4742.3 \\
mpc-congestion   & 18.864 & 16.447 & $-$0.028 & $-$2.151 & 4.597  & 0.7182  & 0.9985 & 4235.1 \\
mpc-hier         & 14.761 & 13.743 & $-$0.085 & $-$1.536 & 2.639  & 0.4124  & 0.9846 & 3733.3 \\
\textbf{mpc-2level} & \textbf{18.898} & \textbf{16.503} & $-$0.030 & $-$2.173 & 4.597 & \textbf{0.7183} & 0.9987 & 4281.9 \\
mpc-hier-route   & 14.779 & 13.811 & $-$0.055 & $-$1.561 & 2.584  & 0.4038  & 0.9848 & 3885.9 \\
\bottomrule
\end{tabular}
\end{table}

\textbf{What drives the gap: admission control, not routing.}
\texttt{mpc-2level} wins outright, but \texttt{mpc-congestion} (predictive
routing only, no full flow-optimization) sits within 0.04 utility of it. A
five-point ISL-rate sweep (\S\ref{sec:sweeps}) found the dedicated
routing-LP's own marginal contribution tops out at $+0.84\%$ utility and is
exactly zero above 1.5\,Mbps ISL. The real, structural $\approx 5$-point gap
is between the no-admission-control family (\texttt{mpc},
\texttt{mpc-congestion}, \texttt{mpc-2level}, $\approx 18.8$--$18.9$) and the
admission-control family (\texttt{mpc-hier}, \texttt{mpc-hier-route},
$\approx 14.8$) — whether to drop a task early matters far more than how
cleverly the kept tasks are routed.

\textbf{Route-freezing does not survive a LEO fabric.} A route frozen for a
10-minute macro-epoch (\texttt{mpc-hier-route}) is usable at its
actually-executed slot only 7.7\% of the time, because ground-station
contact windows (227--265\,s, Table~\ref{tab:constellation}) are shorter
than any freeze period worth having. This is a structural property of the
constellation geometry, not a tuning failure.

\subsection{Solve-time and energy efficiency}

\texttt{mpc-hier}/\texttt{mpc-hier-route} solve $3$--$4\times$ faster per
call (mean 43.5--43.6\,ms) than the flat/two-level MPCs (123.9--168.5\,ms
mean, up to 3{,}088\,ms max for \texttt{mpc-2level}), at the cost of the
utility gap above. \texttt{mpc-2level} solves 365 times over the run
(vs.\ 113--121 for the other MPC variants), for 45.2\,s total solve time
vs.\ 20.4\,s for flat \texttt{mpc} — recall from \S\ref{sec:setup} that these
are individual-solve statistics summed, not simulated wall-clock time.

%   \section{Decision Distributions}
\label{sec:decisions}

This section reports the compression-depth chosen per task, per scheduler,
from the same fabric-limited run (seed 42) as
Tables~\ref{tab:scheduler-comparison}--\ref{tab:utility-components}, so
utility, delivery, and decision-distribution results are directly
comparable. Distributions were computed by aggregating the \texttt{depth}
column of each \texttt{task\_outcomes\_*.csv} (64 tasks per scheduler);
Table~\ref{tab:depth-dist} reports the count of tasks assigned each depth,
plus tasks dropped or rejected before a depth was ever assigned.

\begin{table}[h]
\centering
\caption{Depth-selection distribution per scheduler (of 64 tasks), fabric-limited scenario, seed 42.}
\label{tab:depth-dist}
\begin{tabular}{lrrrrr}
\toprule
Scheduler & $D{=}2$ & $D{=}4$ & $D{=}8$ & $D{=}16$ & Dropped/rejected \\
\midrule
greedy-fixed-2   & 64 & 0  & 0  & 0  & 0 \\
greedy-fixed-4   & 0  & 64 & 0  & 0  & 0 \\
greedy-fixed-8   & 0  & 0  & 64 & 0  & 0 \\
greedy-fixed-16  & 0  & 0  & 0  & 64 & 0 \\
greedy-adaptive  & 0  & 0  & 0  & 64 & 0 \\
mpc              & 0  & 0  & 2  & 62 & 0 \\
mpc-congestion   & 0  & 0  & 1  & 63 & 0 \\
mpc-hier         & 1  & 11 & 35 & 17 & 39 \\
mpc-hier-route   & 1  & 12 & 37 & 14 & 40 \\
mpc-2level       & 0  & 0  & 0  & 64 & 0 \\
\bottomrule
\end{tabular}
\end{table}

\textbf{Fixed-depth greedy baselines behave exactly as specified} (sanity
check: 64/64 tasks at the configured depth for every \texttt{greedy-fixed-*}
row). \textbf{greedy-adaptive converges to depth 16 for every task in this
run} — confirming numerically why its row is identical to
\texttt{greedy-fixed-16} in Tables~\ref{tab:scheduler-comparison}--\ref{tab:utility-components}:
under this scenario's contention level, the adaptive policy's queue-length
trigger never fires away from the maximum depth. \textbf{The no-admission-control
MPC family (mpc, mpc-congestion, mpc-2level) is overwhelmingly depth-16},
using a lower depth only rarely (mpc: 2/64; mpc-congestion: 1/64) and never
dropping a task. \textbf{The admission-control family (mpc-hier,
mpc-hier-route) shows real depth diversity} across all four levels, but this
comes with substantial admission/drop activity (39--40 of 64 tasks dropped
or rejected before delivery) — the mechanism described in
\S\ref{sec:hier} trading deliverability for smaller, faster-solved problems.
This is the same finding stated at the utility level in
\S\ref{sec:utility-results} (\S\ref{sec:formulation} \S\ref{sec:hier}):
admission control, not depth or routing sophistication, is what
distinguishes the two scheduler families.

\pending{This table is reported for one scenario (fabric-limited) and one
seed. Repeating it for the gbs-limited scenario, and for the satellite-/GBS-count
sweep proposed in \S\ref{sec:sweeps}, would show whether the depth-diversity
pattern in the admission-control family is scenario-specific or general.}

%   \section{Sensitivity to Link Rate, Constellation, and Ground-Segment Size}
\label{sec:sweeps}

\subsection{GBS aggregate rate sweep: needs to be rerun}

\texttt{oec\_sim/sweep.py} committed a grid of 6 GBS aggregate rates
$\times$ 5 seeds $\times$ 6 schedulers (180 rows,
\texttt{hypatia\_sim/oec\_scenario/sweep/results.csv}). \pending{This sweep
was generated in the same commit that unified the utility function, but
\emph{before} the real downstream mIoU quality table was wired in one
commit later -- its own commit message states the quality table was ``not
yet run on the server ... provisional until then'' at that point, and the
file has not been regenerated since. Its utility values ($\approx$57--64)
are on the placeholder-quality scale, not the real-quality scale
($\approx$6--19) used everywhere else in this paper. \textbf{This sweep
must be rerun against the current \texttt{quality\_table.json} before its
exact numbers are reported as final} -- flagging this rather than
transcribing stale numbers.} The qualitative shape of the earlier
(placeholder-quality) run is: the best \emph{fixed} depth climbs
$2\to4\to8\to16$ as the GBS rate rises, and \texttt{mpc} beat the best fixed
depth in 5 of 6 rate regimes, tying only at the highest rate tested (no
scarcity, "always max depth" already near-optimal) -- worth confirming
still holds once rerun with real quality.

\subsection{ISL rate sweep}

A parallel 5-point ISL-rate sweep (0.5--2.0\,Mbps, 3 seeds,
\texttt{results\_routing.csv}) isolates the routing MPC's own marginal
contribution: at most $+0.84\%$ utility, and exactly $0\%$ above 1.5\,Mbps
ISL (converges to flat MPC once there is nothing to route around). This
sweep has the identical staleness issue as \S5.1 (placeholder quality
table) and needs the same rerun before its exact numbers are final.

\subsection{Proposed: satellite-count and GBS-count sweeps}

Both axes are existing configuration knobs, not new code:
\texttt{config.CONSTELLATIONS} already defines four named presets varying
satellite count (\texttt{kuiper-630}: 1{,}156 sats; \texttt{starlink-550}:
1{,}584 sats; \texttt{telesat-1015}: 351 sats; \texttt{small-walker}: 6
sats, a topology-validation toy scenario not suited for a results sweep),
and \texttt{config.GROUND\_STATIONS} is a plain list whose length sets
$|\mathcal{G}|$. A satellite-count sweep can therefore run the existing
harness (\texttt{sweep.py}) across the three real constellation presets (or
intermediate synthetic $(\text{planes}, \text{sats/plane})$ tuples added to
\texttt{CONSTELLATIONS}); a GBS-count sweep can add/remove entries from
\texttt{GROUND\_STATIONS} (e.g.\ a 2-, 4-, 8-station subset of Hypatia's
top-100 city list) and rerun. \pending{Neither sweep has been run yet --
this is a proposed design, not a result. Both should use the current real
quality table from the outset to avoid the staleness problem in
\S5.1--5.2.}

%   \section{A Learned Scheduling Policy: PPO vs. MPC}
\label{sec:rl}

Every scheduler in \S\ref{sec:utility-results} solves an optimization problem from scratch at every re-plan. Here we try the alternative: train a policy offline so at run time it only has to evaluate a small network instead of solving a fresh MILP.

\todoconfirm{These results are on \textbf{legacy utility scale} (\S\ref{sec:formulation}), not unified -- the trained policy (\texttt{ppo\_oec.zip}) optimized the old single-factor reward and hasn't been retrained under \texttt{--utility unified}. Don't compare the numbers here against Tables~\ref{tab:scheduler-comparison}--\ref{tab:utility-components}; retraining under the unified reward is future work.}

\subsection{Setup}

The environment (\texttt{oec\_sim/rl\_env.py}) is a Gymnasium environment with a \texttt{Discrete(6)} action space -- commit one of 4 depths, defer, or drop a task -- decided per active task rather than per slot, so the action space doesn't scale with how many tasks are active or how large the constellation is. Reward matches the MILP objective plus a potential-based backlog-shaping term. We trained with PPO for 300k timesteps across load regimes 1.5--2.5\,Mbps and 10 seeds ($\sim$4 minutes on a CPU), then evaluated in-distribution and on a load regime the policy never saw during training.

\subsection{Results}

Table~\ref{tab:rl-vs-mpc} reports utility, delivery, and per-decision compute time for two seeds per regime.

\begin{table}[h]
\centering
\caption{RL (PPO) vs. MPC, legacy utility scale, two seeds per regime.}
\label{tab:rl-vs-mpc}
\begin{tabular}{lrrrrr}
\toprule
Regime & RL utility & RL delivered\% & RL $\mu$s/decision & MPC utility & MPC $\mu$s/decision \\
\midrule
in-dist (2.0\,Mbps)   & 58.13 / 52.31 & 100\% & $\sim$500--1{,}400 & 61.58 / 55.55 & $\sim$145{,}000--243{,}000 \\
OOD-light (4.0\,Mbps) & 58.13 / 52.31 & 100\% & $\sim$500--1{,}400 & 62.44 / 56.14 & $\sim$145{,}000--243{,}000 \\
OOD-heavy (0.75\,Mbps)& 54.76 / 50.83 & 100\% & $\sim$500--700     & 58.92 / 53.27 & $\sim$163{,}000--173{,}000 \\
\bottomrule
\end{tabular}
\end{table}

PPO loses to MPC on utility everywhere tested, by $5$--$8\%$ -- expected, since solving the problem fresh every time beats pattern-matching against training experience. It's also $\sim 250$--$400\times$ faster per decision (microseconds instead of hundreds of milliseconds), which is the actual reason a learned policy might be worth running on power- and compute-constrained flight hardware despite the utility gap.

What we didn't expect: the policy's actions turn out to be provably rate-independent. Its observation includes task properties, active-task count, and mean freshness, but nothing about capacity or congestion, so its decisions are bit-for-bit identical between the 2.0\,Mbps and 4.0\,Mbps regimes at a fixed seed -- visible directly in the table above. It still delivers 100\% of images even at the most constrained regime tested (0.75\,Mbps, where \texttt{greedy-fixed-8} collapses to a fraction of its unconstrained utility), but that's because it happened to learn conservative depth choices, not because it senses congestion. Accidental robustness, not adaptive robustness. The obvious next step is adding a residual-bandwidth feature to the observation, which the original design called for but the shipped version dropped to keep the vector smaller.

Two comparison-fairness bugs in the RL evaluation path were caught and fixed before these numbers were produced: the environment had been crediting RL deliveries with no propagation delay while every other scheduler charged the real route delay (measured effect: $0.0017\%$ utility, negligible against a 30-second slot and 300-second freshness constant -- a correctness fix, not a result change), and the RL scheduler's own utility accounting bypassed the shared \texttt{utility.run\_utility} path, which would have silently dropped the fairness term under \texttt{--utility unified}. Both now go through the same accounting as every other scheduler here.

%   \section{Conclusion}
\label{sec:conclusion}

This paper studied on-board compression depth and downlink scheduling for
Orbital Edge Computing as one coupled decision problem, and found that the
choice of \emph{quality signal} used to reward that decision matters more
than any single scheduler design choice tested. Grounding utility in
measured downstream segmentation performance rather than pixel fidelity
changed MPC's measured advantage over the best fixed-depth baseline from a
marginal $1.2\%$ to a substantial $32.2\%$, because pixel fidelity is
nearly flat across the compression-depth range studied here while
task-relevant quality is not ($+12\%$ vs.\ $+247\%$). A scheduler is only
as good as the signal it is scored against.

With that signal in place, a single unified utility function reconciled
five previously disagreeing scoring definitions, and a ten-scheduler
comparison under a deliberately routing-sensitive capacity regime showed
that the strongest variant (\texttt{mpc-2level}, a routing MPC and a depth
MPC solved as peers) closes to within $0.5\%$ of the offline optimality
bound. The more consequential finding, however, was architectural rather
than numerical: the $\approx 5$-point utility gap that separates the
strongest and weakest MPC variants tracks \emph{admission control}
(whether a scheduler explicitly drops an overwhelmed task rather than let
it rot in queue), not routing sophistication — the dedicated routing
optimizer's own marginal contribution tops out at $+0.84\%$ utility and
disappears entirely once link rates are not scarce. A related structural
result is that freezing a concrete route for longer than a few minutes
does not survive this constellation's short ground-contact windows,
regardless of how that freeze period is tuned. Finally, a learned (PPO)
policy traded a consistent $5$--$8\%$ utility gap for a
$250$--$400\times$ reduction in per-decision compute time, with an
unanticipated caveat: without a congestion feature in its observation, its
apparent robustness to link scarcity is accidental rather than adaptive.

\paragraph{What remains open.} Several results in this paper are
explicitly flagged, not silently finalized: the uncompressed-reference
mIoU used to anchor the quality signal has not been verified against
FLAIR-1's official held-out test split (\S\ref{sec:downstream}); the scope
of \texttt{mpc-hier} relative to Xuanhao's request to drop certain prior
results is pending his direct confirmation (\S\ref{sec:utility-results});
and both link-rate sensitivity sweeps in \S\ref{sec:sweeps} were generated
against a placeholder quality table and need to be rerun against the real,
measured one before their exact numbers are reported as final. The
proposed satellite-count and ground-station-count sweeps
(\S\ref{sec:sweeps}) and a retraining of the PPO baseline under the
unified reward (\S\ref{sec:rl}) are natural next steps that require no new
code, only additional runs. None of these open items affect the paper's
central conclusion — that a scheduler's measured value depends first on
what it is being asked to optimize, and only second on how cleverly it
optimizes it — but resolving them would sharpen every number built on top
of that conclusion.

% \bibliographystyle{unsrt}
% \bibliography{reference}

\end{document}
